\section{The monodromy group and its period domain}\label{fam:section:F}

We use the clockwise root deck generators and the lattice basis specified
on the first page.  In this section a group acts properly discontinuously when, for every
compact subset $K$ of its space, only finitely many group elements $g$ satisfy
$gK\cap K\ne\varnothing$.  This definition permits finite stabilizers.  We
write $e=\widehat\delta$, $E=e\otimes\gamma$, and
\[
 J_2=\begin{pmatrix}0&1\\-1&0\end{pmatrix},\qquad
 \Omega=\begin{pmatrix}
 0&0&0&6\\0&0&1&0\\0&-1&0&0\\-6&0&0&0
 \end{pmatrix}.
\]
Thus the alternating form is
$\psi=u\wedge w+6\gamma\wedge\delta$, an integral alternating form of
type $(1,6)$ on $\Lambda$.  We use $\Omega$ for its matrix and retain $J$ for
a real complex structure.  The domain $\mathcal D$ consists of the real
endomorphisms $J$ of $\Lambda_{\mathbb R}$ satisfying
$J^2=-I$, $J^{\mathsf t}\Omega J=\Omega$, and
$\operatorname{sign}(\Omega J)=(2,2)$.  It is the stated
$\operatorname{Sp}_4(\mathbb R)/U(1,1)$, with its complex structure as an
open subset of the Lagrangian Grassmannian.  An integral symplectic matrix
acts by $J\mapsto gJg^{-1}$.

\subsection{An exact arithmetic description}\label{fam:section:F.1}

\begin{theorem}\label{fam:src:F.1.1}\label{fam:st:F.1.1}\label{cdp:st:D.1.14}
Let $P_{e,\mathbb Z}$ be the integral symplectic group fixing $e$ as a
vector, and let $P_{\mathbb Z}$ be the integral symplectic group preserving
$\mathbb Ze\subset\ker\gamma\subset\Lambda$.  Every element of
$P_{e,\mathbb Z}$ has a unique expression
\begin{equation}\label{fam:eq:F-g}
 g(B,x,t)=
 \begin{pmatrix}
 1&0&0\\6x&B&0\\t&-x^{\mathsf t}J_2B&1
 \end{pmatrix},\qquad
 B\in\operatorname{SL}_2(\mathbb Z),\quad x\in\mathbb Z^2,\quad t\in\mathbb Z,
\end{equation}
where the block sizes are $1,2,1$.  Put
\[
 S=\begin{pmatrix}0&-1\\1&0\end{pmatrix},\quad
 T=\begin{pmatrix}1&1\\0&1\end{pmatrix},\quad
 a=\begin{pmatrix}0&1\\-1&-1\end{pmatrix}.
\]
There is a homomorphism
$\alpha:\operatorname{SL}_2(\mathbb Z)\to\mathbb Z/12$ determined by
$\alpha(S)=9$ and $\alpha(T)=1$; in particular $\alpha(a)=4$.
With $Q(x)=x_1+x_2+x_1x_2$, the actual monodromy group is
\begin{equation}\label{fam:eq:F-Gamma}
 \Gamma=\left\{g(B,x,t):t+6Q(x)-\alpha(B)=0\ \text{in }\mathbb Z/12\right\}.
\end{equation}
Consequently $[P_{e,\mathbb Z}:\Gamma]=12$ and
$P_{\mathbb Z}=\{\pm I\}\times P_{e,\mathbb Z}$, so
$[P_{\mathbb Z}:\Gamma]=24$.  Its image in the Levi quotient is exactly
$\operatorname{SL}_2(\mathbb Z)$, with scalar on $e$ equal to $1$.
\end{theorem}
\Uses{the hypotheses, notation and definitions stated here; \cref{fam:eq:F-law}; \cref{fam:eq:F-SLpresentation}}


\begin{proof}
The matrices in the stated problem, in the fixed basis, are
\[
 A_1=\begin{pmatrix}1&0&0&0\\6&0&1&0\\-6&-1&-1&0\\-2&1&0&1\end{pmatrix},
 \qquad
 A_2=\begin{pmatrix}1&0&0&0\\0&0&-1&0\\-6&1&0&0\\3&0&1&1\end{pmatrix}.
\]
Multiplication gives $A_1^3=A_2^4=I$ and
$A_j^{\mathsf t}\Omega A_j=\Omega$.  In particular
$A_1=g(a,(1,-1)^{\mathsf t},-2)$ and
$A_2=g(S,(0,-1)^{\mathsf t},3)$.  Their images $a,S$ generate
$\operatorname{SL}_2(\mathbb Z)$ because $Sa=T$ and $S,T$ perform the
Euclidean algorithm on primitive integral columns.

A matrix fixing $e$ and preserving $\Omega$ preserves
$e^{\perp}=\ker\gamma$.  Its first row is $(1,0,0,0)$, its last column is
$e$, its middle block $B$ is symplectic on
$e^{\perp}/\mathbb Re$, and its remaining middle column, say $v$, satisfies
$6r=-v^{\mathsf t}J_2B$.  Integrality of $r$ and unimodularity of $B$
give $v=6x$ with $x\in\mathbb Z^2$.  There is no further condition on
the lower left entry $t$.  This proves \eqref{fam:eq:F-g}.  Direct multiplication
gives
\begin{equation}\label{fam:eq:F-law}
 g(B,x,t)g(C,y,s)
 =g(BC,x+By,t+s-6x^{\mathsf t}J_2By).
\end{equation}

For completeness, the presentation of the Levi group needed here follows
from an elementary action on the real projective line.  The images of $a$
and $S$ in $\operatorname{PSL}_2(\mathbb Z)$ have orders $3$ and $2$.
On positive real numbers, $a(z)=-1/(z+1)$ has image $(-1,0)$ and
$a^2(z)=-(z+1)/z$ has image $(-\infty,-1)$; $S$ carries negative real
numbers to positive real numbers.  A cyclically reduced alternating word
starting with a nonidentity power of $a$ and ending with $S$ therefore
carries $(-\infty,0)$ into one of these two proper subintervals.  It cannot
be the identity.  Every reduced word is conjugate to such a word or to a
nonidentity element of one factor.  Together with the Euclidean algorithm,
this proves the presentation
$\operatorname{PSL}_2(\mathbb Z)=\langle\bar a,\bar S\mid
\bar a^3=\bar S^2=1\rangle$.  It follows that
\begin{equation}\label{fam:eq:F-SLpresentation}
 \operatorname{SL}_2(\mathbb Z)=
 \langle a,S\mid a^3=S^4=1,\ [a,S^2]=1\rangle.
\end{equation}
Indeed the group on the right modulo its central element $S^2$ is the
just-proved free product, and $S^2$ maps to the nonidentity matrix $-I$.
Its kernel in the displayed matrix group is consequently trivial.
The values $4,9$ respect all these relations, proving the assertion about
$\alpha$; $T=Sa$ has value $1$.

The quadratic function $Q$ modulo $2$ is zero at $(0,0)$ and one at all
three other points of $\mathbb F_2^2$.  Hence $Q(Bx)=Q(x)\pmod2$ for every
$B\in\operatorname{SL}_2(\mathbb Z)$.  Also
\[
 Q(x+y)-Q(x)-Q(y)\equiv x^{\mathsf t}J_2y\pmod2.
\]
Equation \eqref{fam:eq:F-law} now shows that
$h(g(B,x,t))=t+6Q(x)\pmod{12}$ is a homomorphism.
The homomorphism $\chi=h-\alpha\circ B$ is surjective, since
$\chi(g(I,0,1))=1$.  Its values on $A_1,A_2$ are zero.

We verify equality with its kernel; containment alone would not determine
the group.  Let $N_{\mathbb Z}$ be the kernel of the map to the middle
block.  The kernel of
$\langle g_1,g_2\mid g_1^3=g_2^4=1\rangle\to
\operatorname{SL}_2(\mathbb Z)$ is normally generated by
$[g_1,g_2^2]$, by \eqref{fam:eq:F-SLpresentation}.  Its image in $\Gamma$ is
therefore normally generated by
\[
 K=[A_1,A_2^2]=g(I,(0,-1)^{\mathsf t},6).
\]
Put
\[
 U=A_2KA_2^{-1}=g(I,(1,0)^{\mathsf t},-6),\quad
 V=K^{-1}=g(I,(0,1)^{\mathsf t},-6),\quad Z=I+12E.
\]
Equation \eqref{fam:eq:F-law} gives $[U,V]=Z^{-1}$.  Thus $U,V,Z$ generate
\begin{equation}\label{fam:eq:F-H}
 H=\{g(I,x,t):t\in6\mathbb Z,\quad t/6\equiv Q(x)\pmod2\}.
\end{equation}
More explicitly,
$U^rV^sZ^k=g(I,(r,s)^{\mathsf t},-6(r+s+rs)+12k)$, so these expressions
give every member of the set.  Conjugation by $g(B,x,t)$ sends
$g(I,y,s)$ to $g(I,By,s-12x^{\mathsf t}J_2By)$; the invariance of $Q$
modulo $2$ proves that $H$ is normal in $P_{e,\mathbb Z}$.  The normal
closure of $K$ is contained in $H$ and contains $U,V,Z$, so it is exactly
$H$.  We have proved $\Gamma\cap N_{\mathbb Z}=H$.  Both $\Gamma$ and
$\ker\chi$ surject onto the same Levi group and have intersection $H$
with $N_{\mathbb Z}$.  Given $g\in\ker\chi$, choose $\gamma\in\Gamma$
with the same middle block.  Then $g\gamma^{-1}\in H$, and $g\in\Gamma$.
This proves \eqref{fam:eq:F-Gamma} and the index $12$.

Finally an integral automorphism preserving the primitive line
$\mathbb Ze$ sends $e$ to $e$ or $-e$.  Multiplication by the central
matrix $-I$ converts the latter case to the former.  The intersection of
$\{\pm I\}$ with $P_{e,\mathbb Z}$ is trivial.  This proves the direct
product and the index $24$.
\end{proof}

\begin{proposition}\label{fam:src:F.1.2}\label{fam:st:F.1.2}
The group in \eqref{fam:eq:F-Gamma} has center
$\langle Z\rangle\cong\mathbb Z$ and is the following explicit central
extension of $\mathbb Z^2\rtimes\operatorname{SL}_2(\mathbb Z)$ by
$\mathbb Z$.  Choose the representative $a(B)\in\{0,\ldots,11\}$ of
$\alpha(B)$.  Its underlying set consists of triples $(B,x,k)$, with
product
\[
 (B,x,k)(C,y,l)=(BC,x+By,k+l+\kappa((B,x),(C,y))),
\]
where
\begin{align*}
12\kappa((B,x),(C,y))={}&a(B)+a(C)-a(BC)-6Q(x)-6Q(y)\\
 &-6x^{\mathsf t}J_2By+6Q(x+By).
\end{align*}
Its real and rational Zariski closures are respectively the full real
and rational symplectic stabilizers of the vector $e$.  They have
dimension $6$, with three-dimensional Heisenberg unipotent radical and
Levi quotient $\operatorname{SL}_2$.
\end{proposition}
\Uses{the hypotheses, notation and definitions stated here; \cref{fam:eq:F-Gamma}; \cref{fam:eq:F-law}}


\begin{proof}
The numerator defining $\kappa$ is divisible by $12$, by the two
congruences proved above.  The bijection
\[
 (B,x,k)\longmapsto g(B,x,a(B)-6Q(x)+12k)
\]
intertwines the displayed multiplication with matrix multiplication by
\eqref{fam:eq:F-law}.  This proves associativity as well as the claimed abstract
description, without leaving an unspecified extension class.  The kernel
of the homomorphism to $\mathbb Z^2\rtimes\operatorname{SL}_2(\mathbb Z)$
is $\langle Z\rangle$.  That affine group has trivial center: commuting
with all translations forces its linear part to be $I$, and commuting
with $S,T$ forces its translation part to vanish.  Thus the center of
$\Gamma$ is exactly the stated kernel.

The logarithms $U-I,V-I,Z-I$ are linearly independent elements of the
three-dimensional Heisenberg Lie algebra.  The square of each is zero.
A polynomial vanishing on all integral powers $I+n(U-I)$ vanishes on
the whole one-parameter line; the same holds for $V$ and $Z$.  Hence the
Zariski closure contains the full unipotent radical.  Its Levi image
contains $S,T$; the powers of $T$ and of $STS^{-1}$ similarly supply the
two opposite unipotent subgroups of $\operatorname{SL}_2$.  Their Lie
algebras and their bracket span $\mathfrak{sl}_2$.  They generate the
connected algebraic group $\operatorname{SL}_2$.  A subgroup containing
the entire unipotent kernel and surjecting onto this Levi quotient is
the entire vector stabilizer.  Its dimension is $3+3=6$.  This reasoning
uses equations over $\mathbb Q$, so proves the rational assertion as
well as the real assertion.
\end{proof}

\subsection{Properness and the maximal open set}\label{fam:section:F.2}

Define the real-analytic function
\[
 q_e:\mathcal D\longrightarrow\mathbb R,\qquad q_e(J)=\psi(e,Je),
 \qquad\mathcal O=\{J:q_e(J)\ne0\}.
\]
The function $q_e$ is invariant under $P_{\mathbb Z}$: replacing $e$ by
$-e$ changes both of its arguments.  Write $\mathcal O_+$ and
$\mathcal O_-$ for its two sign regions.

\begin{proposition}[counterexample]\label{fam:src:F.2.1}\label{fam:st:F.2.1}
The assertion that $\Gamma$ acts properly discontinuously on all of
$\mathcal D$ is false.  The following point of $\mathcal D$ has an
explicit infinite stabilizer:
\[
 J_0=\begin{pmatrix}
 0&0&-1&0\\0&0&0&-6\\1&0&0&0\\0&1/6&0&0
 \end{pmatrix},\qquad
 R=\begin{pmatrix}
 1&0&0&0\\0&1&72&0\\0&0&1&0\\-12&0&0&1
 \end{pmatrix}\in\Gamma.
\]
Every power of $R$ fixes $J_0$, and all these powers are distinct.
\end{proposition}
\Uses{the hypotheses, notation and definitions stated here; \cref{fam:eq:F-Gamma}}


\begin{proof}
Direct multiplication gives $J_0^2=-I$ and
$J_0^{\mathsf t}\Omega J_0=\Omega$.  Moreover
\[
 \Omega J_0=\begin{pmatrix}
 0&1&0&0\\1&0&0&0\\0&0&0&6\\0&0&6&0
 \end{pmatrix}
\]
has eigenvalues $1,-1,6,-6$, so $J_0\in\mathcal D$ and $q_e(J_0)=0$.
The matrix $R$ equals $g(T^{72},0,-12)$ and belongs to \eqref{fam:eq:F-Gamma},
because $\alpha(T^{72})=0\pmod{12}$.  Its commutator with $J_0$ is zero.
Writing $R=I+N$ gives $N\ne0$ and $N^2=0$, so $R^m=I+mN$ and the
powers are distinct.  The compact singleton $\{J_0\}$ therefore violates
the stated definition of proper discontinuity.  These identities are
checked both with SymPy and with a separate rational matrix
implementation in \texttt{code/f\_exact.py}.
\end{proof}

\begin{theorem}\label{fam:src:F.2.2}\label{fam:st:F.2.2}
The set $\mathcal O$ is the largest $\Gamma$-invariant open subset of
$\mathcal D$ on which the action is properly discontinuous.  It is also
the largest such open subset for $P_{\mathbb Z}$.  Each sign region is
a complex manifold in the analytic topology, and its quotient by either
discrete group is a Hausdorff complex analytic orbifold.
\end{theorem}
\Uses{the hypotheses, notation and definitions stated here; \cref{fam:eq:F-g}}


\begin{proof}
Let $P_{\mathbb R}$ be the real symplectic group preserving the real line
$\mathbb Re$.  We first prove properness on each sign region.  For a fixed
$J$, the form
$h_J(v,w)=\psi(v,Jw)+i\psi(v,w)$ is a nondegenerate Hermitian form of
signature $(1,1)$ on the complex vector space $(\Lambda_{\mathbb R},J)$,
with the convention conjugate linear in the first entry.  When
$q_e(J)\ne0$, the complex line $\mathbb Ce$ is nondegenerate and its
Hermitian orthogonal complement is a complex line of the opposite sign.
Normalize $e$ by a positive real scalar and choose a unit vector in this
complement.  Comparing these two Hermitian bases for two complex
structures with the same sign gives a real symplectic map preserving
$\mathbb Re$ and conjugating one complex structure to the other.  Thus
$P_{\mathbb R}$ is transitive on each sign region.

An element of the stabilizer of $J$ inside $P_{\mathbb R}$ is complex
linear and sends $e$ to $r e$ for a real nonzero $r$.  Preservation of
$q_e(J)$ gives $r^2=1$.  On the complementary complex line it acts by an
element of $U(1)$.  The stabilizer is consequently the compact group
$\{\pm1\}\times U(1)$.  Each sign region is therefore a homogeneous
space of $P_{\mathbb R}$ with compact stabilizer $L$.  Local sections of
$P_{\mathbb R}\to P_{\mathbb R}/L$, and a finite cover of a compact
subset $K$, give a compact subset $C$ of $P_{\mathbb R}$ mapping onto
$K$.  Every element carrying a point of $K$ back into $K$ belongs to
$CLC^{-1}$, which is compact.  This proves properness.  Restricting a
proper action to a closed discrete subgroup gives
the required finite-intersection condition.  Both $\Gamma$ and
$P_{\mathbb Z}$ are closed discrete matrix groups.  The two sign regions
are disjoint invariant open-and-closed subsets of $\mathcal O$, so the
same finite-intersection condition holds on their union.  Proper actions
by discrete groups have finite stabilizers and Hausdorff quotients; their
holomorphic local finite-group quotient charts supply the asserted
analytic orbifolds.

It remains to prove maximality, which requires more than one bad point.
The real vector stabilizer $P_{e,\mathbb R}$ acts transitively on
$\mathcal D_0=\{J:q_e(J)=0\}$.  Here is a direct comparison map.  In the
Hermitian space just described, $e$ is a nonzero null vector.  Choose
$v$ with $h_J(e,v)=1$, and replace $v$ by
$v-\tfrac12h_J(v,v)e$.  This preserves $h_J(e,v)=1$ and makes $v$ null.
The ordered complex basis $(e,v)$ has Gram matrix
$\bigl(\begin{smallmatrix}0&1\\1&0\end{smallmatrix}\bigr)$.
The complex-linear map between two such bases fixes $e$ and preserves
$h$, hence its imaginary part $\psi$.  As a real map it belongs to
$P_{e,\mathbb R}$ and conjugates the two complex structures.

The subgroup $P_{e,\mathbb Q}$ is dense in $P_{e,\mathbb R}$.
For the unipotent coordinates in \eqref{fam:eq:F-g} this is density of
$\mathbb Q^3$ in $\mathbb R^3$.  For $\operatorname{SL}_2$ one may
approximate $a,b,c$ by rational numbers in the chart $a\ne0$ and set
$d=(1+bc)/a$; translates of that chart cover the group.  Hence
\[
 \{rJ_0r^{-1}:r\in P_{e,\mathbb Q}\}
\]
is dense in $\mathcal D_0$.
Put $R_0=g(T^6,0,-1)$, so $R_0^{12}=R$ and
$(R_0-I)^2=0$.  At $rJ_0r^{-1}$ the nonzero rational nilpotent
$L=r(R_0-I)r^{-1}$ commutes with the complex structure and satisfies
$L^2=0$.  Choose a positive integer $d$ clearing its denominators.
The matrix $I+dL$ and its inverse $I-dL$ are integral, symplectic, and
fix $e$; this follows either by conjugating the identity
$(I+s(R_0-I))^{\mathsf t}\Omega(I+s(R_0-I))=\Omega$ or by expanding it.
The twelfth power of $I+dL$ lies in $\Gamma$, because
$P_{e,\mathbb Z}/\Gamma\cong\mathbb Z/12$.  It has infinite order and
fixes $rJ_0r^{-1}$.  Thus points with infinite $\Gamma$-stabilizer are
dense in $\mathcal D_0$.

Every open neighborhood in $\mathcal D$ of any point of $\mathcal D_0$
therefore contains a point with infinite stabilizer.  An open set on
which the action is proper cannot meet $\mathcal D_0$.  This proves the
largest-open-set assertion for $\Gamma$.  The same points have infinite
$P_{\mathbb Z}$-stabilizer, so it proves that assertion for the full
parabolic as well.
\end{proof}

\begin{proposition}\label{fam:src:F.2.3}\label{fam:st:F.2.3}
The positive sign region admits the global period chart
\[
 \mathcal O_+=\left\{
 (\tau,\mu,\beta)\in\mathbb C^3:
 \operatorname{Im}\tau>0,\quad
 \operatorname{Im}\beta-\frac{6(\operatorname{Im}\mu)^2}
 {\operatorname{Im}\tau}<0\right\},
\]
represented by the matrix $\Pi$ in the stated problem.  In this chart,
\begin{equation}\label{fam:eq:F-qD}
 q_e(J)=-\frac6D.
\end{equation}
Every admissible period $\Pi_c(z)$ belongs to $\mathcal O_+$, for every
admissible $c$.  Thus this single maximal proper-action open set contains
all the period images in question.
\end{proposition}
\Uses{the hypotheses, notation and definitions stated here; the explicit constructions and calculations in the proof}


\begin{proof}
Let $J\in\mathcal O_+$.  The two vectors $\widehat w,e$ are complex-linearly
independent in $(\Lambda_{\mathbb R},J)$.  Otherwise
$\widehat w=r e+sJe$ for real $r,s$.  Pairing with $e$ gives
$0=sq_e(J)$; hence $s=0$, contradicting the real linear independence
of $\widehat w,e$.  There is consequently a unique complex-linear
isomorphism $(\Lambda_{\mathbb R},J)\to\mathbb C^2$ sending these two
vectors to $(1,0)$ and $(0,1)$.  The corresponding period matrix has
right block $I$.  The equation
$\Pi\Omega^{-1}\Pi^{\mathsf t}=0$ says that its upper left entry is
six times its lower middle entry, giving exactly the displayed matrix.
This construction is the ordinary holomorphic Grassmannian chart,
so its coordinates and inverse are holomorphic.

In the real symplectic basis
$(\widehat u,\widehat\gamma/6,\widehat w,e)$ the normalized symmetric
period matrix is
$\bigl(\begin{smallmatrix}\tau&\mu\\\mu&\beta/6\end{smallmatrix}\bigr)$.
Its imaginary part has signature $(1,1)$, equivalently
$6(\operatorname{Im}\mu)^2-
\operatorname{Im}\tau\operatorname{Im}\beta>0$.
Solving $\Pi(Je)=i\Pi e=(0,i)$ in real coefficients gives
\[
 \gamma(Je)=
 \frac{\operatorname{Im}\tau}
 {\operatorname{Im}\tau\operatorname{Im}\beta-
 6(\operatorname{Im}\mu)^2}.
\]
Since $\psi(e,v)=-6\gamma(v)$, positivity of $q_e$ now gives
$\operatorname{Im}\tau>0$ and $D<0$, and proves \eqref{fam:eq:F-qD}.
Conversely these inequalities make the imaginary period matrix
invertible with the stated signature and give $q_e>0$.  The admissibility
inequality in the stated problem is $D_c<0$ everywhere, which proves the last
assertion by this exact comparison map.
\end{proof}

\subsection{The curve in the proper quotient}\label{fam:section:F.3}

\begin{theorem}\label{fam:src:F.3.1}\label{fam:st:F.3.1}
For every admissible $c$, the marked period map on the covering of
$B^\circ$ is a holomorphic immersion into $\mathcal O_+$.  The induced
map
\[
 p_c:B^\circ\longrightarrow\Gamma\backslash\mathcal O_+
\]
is injective and a holomorphic immersion.  It extends over $t=0,1$ to a
proper holomorphic embedding
\[
 \bar p_c:\mathbb C=B\setminus\{\infty\}
 \lhook\joinrel\longrightarrow\Gamma\backslash\mathcal O_+.
\]
The closure of $p_c(B^\circ)$ in this quotient is exactly
$\bar p_c(\mathbb C)$.
\end{theorem}
\Uses{the hypotheses, notation and definitions stated here; \cref{fam:eq:F-j}; \cref{fam:src:F.2.2}}


\begin{proof}
For $B=\bigl(\begin{smallmatrix}a&b\\c&d\end{smallmatrix}\bigr)$, changing
the marking by $g(B,x,t)$ replaces $\Pi$ by $\Pi g(B,x,t)^{-1}$ and then
renormalizes its right block.  On the quotient by the complex line of
$e$, the two periods are the row $(\tau,1)$.  Right multiplication by
$B^{-1}$ sends it to $(d\tau-c,a-b\tau)$.  Hence
\begin{equation}\label{fam:eq:F-tauaction}
 \tau\longmapsto\frac{d\tau-c}{a-b\tau}.
\end{equation}
The denominator never vanishes in the upper half plane.  The function
$j(\tau)$ is therefore a holomorphic invariant of this action and
descends to a holomorphic function $j_\Gamma$ on the proper quotient.
Here $j$ is the elliptic invariant, normalized by
$j(E)=1728\,4A^3/(4A^3+27B^2)$ for $y^2=x^3+Ax+B$; changing a lattice
basis preserves the isomorphism of elliptic curves and hence this
function.

We derive the required value in the present family.  The elliptic curve
whose period is $\tau$ has the equation stipulated in the stated problem,
\[
 y^2=x^3-3t^3(t-1)x+2t^4(t-1)^2.
\]
Its discriminant is $1728t^8(t-1)^3$ and direct cancellation gives
\begin{equation}\label{fam:eq:F-j}
 j(\tau)=1728t\qquad(t\in B^\circ).
\end{equation}
Equivalently this is the defining lift of the elliptic modular map used
for $\tau$ in the period construction.  The computation is reproduced
exactly in \texttt{f\_exact.py}.  In a local base coordinate $t$, the
derivative of the composite $j\circ\tau$ equals $1728$.  Thus the
derivative of the marked period map is nonzero.  The same identity after
passage to the quotient shows both injectivity and injectivity of the
tangent map of $p_c$: $j_\Gamma/1728$ is a holomorphic left inverse.

At the two elliptic points, the period functions are holomorphic in the
cyclic covering coordinates $s_1,s_2$, with $t=s_1^3$ and
$t-1=s_2^4$ after choosing local coordinates as in the stated problem.  Their
values are nondegenerate because $D_c<0$ there.  The maps from these
discs to $\mathcal O_+$ are equivariant for $A_1,A_2$.  Composing with
the quotient makes them invariant, so they descend to holomorphic maps
from the base discs.  This is the finite-quotient comparison map: a
holomorphic invariant germ on a disc is a convergent series in $s_j^{m_j}$,
which proves descent on local function rings.  The extensions agree
with $p_c$ off their centers and therefore with each other.  Identity
\eqref{fam:eq:F-j} extends by holomorphy, so the resulting $\bar p_c$ still
has holomorphic left inverse $j_\Gamma/1728$.

The quotient is Hausdorff by \ref{fam:src:F.2.2}.  For a compact subset $K$ of it,
$\bar p_c^{-1}(K)$ is closed in $\mathbb C$ and lies in the compact set
$j_\Gamma(K)/1728$.  Thus $\bar p_c$ is proper.  A holomorphic section of
a holomorphic map is locally a closed immersion: in local embeddings
into polydiscs, its coordinate along the base is $t$ and the remaining
coordinates are its holomorphic component functions, whose graph
defines the image.  Applying this to the left inverse shows that
$\bar p_c$ is a holomorphic embedding; properness makes its image
closed.  Since $\mathbb C\setminus\{0,1\}$ is dense in $\mathbb C$, its
image has the asserted closure.
\end{proof}

\begin{proposition}\label{fam:src:F.3.2}\label{fam:st:F.3.2}
At the two finite punctures the period map has the holomorphic
finite-cover limits
\[
 \Pi_1=
 \begin{pmatrix}2(2-\rho)&\rho&1&0\\
 \beta_c(z_1)&(2-\rho)/3&0&1\end{pmatrix},\qquad
 \Pi_2=
 \begin{pmatrix}3(1-i)&i&1&0\\
 \beta_c(z_2)&(1-i)/2&0&1\end{pmatrix},
 \quad \rho=e^{\pi i/3}.
\]
Their monodromies are the given order-three and order-four matrices.
At the cusp the curve leaves every compact subset of
$\Gamma\backslash\mathcal O_+$.  On the prescribed cusp branch its
Hodge plane in the compact dual tends to
$\mathbb C\gamma+\mathbb Cu$.  The latter plane has complex dimension
two in common with its conjugate and therefore does not represent a
point of $\mathcal D$.
\end{proposition}
\Uses{the hypotheses, notation and definitions stated here; \cref{fam:src:F.3.1}}


\begin{proof}
The finite values of $\tau,\mu$ follow by evaluating their transformation
laws at the fixed points: $\tau(z_1)=\rho$,
$\mu(z_1)=(2-\rho)/3$, $\tau(z_2)=i$,
$\mu(z_2)=(1-i)/2$.  Substitution in $\Pi$ gives the two matrices.
The finite covering comparison and its descent were proved in \ref{fam:src:F.3.1}.
The unipotent cusp transformation is
$\tau\mapsto\tau-1$, $\mu\mapsto\mu$, $\beta\mapsto\beta+1$.
On the cusp branch, $\operatorname{Im}\tau\to+\infty$,
$\mu\to\mu_0$, and $\beta_c+\tau\to b_0+c$, from the holomorphic
cusp expansions.  Consequently
\[
 D_c=-\operatorname{Im}\tau+O(1),\qquad
 q_e(\Pi_c)\longrightarrow0^+,
 \qquad j_\Gamma(p_c(t))=1728t\longrightarrow\infty.
\]
The last equality alone excludes a convergent subnet in any compact
subset of the proper quotient; the first two identities give the
position relative to the maximal proper-action open set.
The rows of $\Pi_c$, regarded as covectors, are
$6\mu\gamma+\tau u+w$ and $\beta_c\gamma+\mu u+\delta$.
Dividing each by $\tau$ gives limits $u$ and $-\gamma$.
Their exterior product has a nonzero limit, so their span converges
in the Grassmannian to the stated plane.  It is spanned by real
covectors, hence equals its conjugate.  A complex structure has
$H^{1,0}\cap H^{0,1}=0$; here the intersection is the entire
two-dimensional plane.  This explicit comparison proves the assertion
about the compact-dual limit.
\end{proof}

\begin{proposition}\label{fam:src:F.3.3}\label{fam:st:F.3.3}
The determinant used in the workbench is generically nonzero, with
the stated zero at $z_1$:
\[
 \det\dot Z=6\dot\mu^2-\dot\tau\dot\beta,\qquad
 \det\dot Z(z_1)=0.
\]
The immersion established in \ref{fam:src:F.3.1} holds at every point of $B^\circ$
without requiring this determinant to be everywhere nonzero.
\end{proposition}
\Uses{the hypotheses, notation and definitions stated here; \cref{fam:src:F.3.1}; \cref{fam:eq:F-j}}


\begin{proof}
In the cusp coordinate $s$, the expansion
$Z=s\bigl(\begin{smallmatrix}0&1\\-1&0\end{smallmatrix}\bigr)+C(q)$,
$q=e^{2\pi is}$, gives
$\det(dZ/ds)=1+O(q)$, so the determinant is not the zero function.
At $z_1$ the fixed-point transformation law of $\mu$ gives
$\dot\mu=-\dot\tau/3$.  One may obtain this by differentiating
$\mu\circ g_1=(1-\mu)/\tau$, using
$g_1'(z_1)=-\rho$ and $\rho^2-\rho+1=0$.
The law
$\beta\circ g_1=\beta+2-6(1-\mu)^2/\tau$ then gives
$\dot\beta=2\dot\tau/3$.  Substitution yields
$6\dot\tau^2/9-2\dot\tau^2/3=0$.
The derivative of a map from a curve to the three-dimensional period
chart is a vector of three coordinate derivatives.  Identity
\eqref{fam:eq:F-j} proves that this vector is nonzero over $B^\circ$;
the determinant tests the rank of the associated two-by-two
Kodaira--Spencer matrix.  The maps compared here are explicitly the
coordinate derivative and its determinant, and the displayed computation
states exactly what each supplies.
\end{proof}



\section{Exceptional fibres, Neron--Severi classes, and compact subspaces}\label{fam:section:D}
\subsection{Elliptic quotients and the torsion locus}\label{fam:section:D.1}
\begin{proposition}\label{fam:src:D.1.1}\label{fam:st:D.1.1}
The specified linear projection \(\operatorname{pr}_1:\mathbb C^2\to
\mathbb C\) induces a holomorphic group homomorphism
\(F_t\to E_t=\mathbb C/(\mathbb Z\tau+\mathbb Z)\) exactly on
\(6\mu=A\tau+B\), \(A,B\in\mathbb Z\). It is then surjective with
connected kernel
\(\mathbb C/(\mathbb Z+(A\mu-\beta)\mathbb Z)\).
On the smooth base the two values of \(t\) are precisely
\[
 t_\pm=\frac{64(47\pm27\sqrt3)}{11}.
\]
\end{proposition}
\Uses{the hypotheses, notation and definitions stated here; the explicit constructions and calculations in the proof}

\begin{proof}
The images of the four lattice generators under the linear projection
are \(6\mu,\tau,1,0\). This proves descent in both directions.
The integral basis change replacing \(\widehat\gamma\) with
\(\widehat\gamma-A\widehat u-B\widehat w\) changes its period to
\((0,\beta-A\mu)\); the remaining vertical generator is \((0,1)\).
The other two periods map onto the entire target lattice, so the kernel
is connected. Its second modulus has imaginary part
\(A\operatorname{Im}\mu-\operatorname{Im}\beta=-D>0\).

The elliptic model may be scaled by the explicit isomorphism
\((X,Y)\mapsto(2X,2\sqrt2Y)\) to
\[
 Y^2=X^3-\frac34t^3(t-1)X+\frac14t^4(t-1)^2,
 \qquad P=(0,t^2(t-1)/2).
\]
Its sixth division polynomial at \(P\) is
\[
 \psi_6(P)=\frac9{65536}t^{24}(t-1)^9
 (11t^2-6016t+8192).
\]
For a smooth elliptic fibre, \(\psi_n(P)=0\) is exactly the condition
\(nP=O\), since the multiplication rational map has these and only these
poles in its affine coordinate. The two roots displayed above are simple
and avoid \(0,1\). We computed the same equation independently by the
chord-and-tangent addition algorithm, obtaining
\[
 \begin{aligned}
 x(6P)&=\frac{(89t^2+4736t-4096)Q(t)}
 {1296(11t^2-6016t+8192)^2},\\
 Q(t)&=881t^4+4745344t^3-12079104t^2\\
 &\hspace{2em}+24641536t-16777216.
 \end{aligned}
\]
The division-polynomial recurrence and the addition algorithm are in
separate supplied programs; all polynomial operations use exact rational
arithmetic. Thus the denominator has no unverified numerical root or
cancellation. The decimal values \(1.36510950191\) and
\(545.543981407\) are only convenient approximations to the exact roots.
\end{proof}


\paragraph{Cited input.} \label{fam:src:D.1.2}\label{fam:st:D.1.2}The source statement and its expanded derivation are retained in the proof compendium, Paper~I, label \nolinkurl{fam:src:D.1.2}.\par

\Uses{the hypotheses, notation and definitions stated here; \cref{app:rederive:fam-src-D.1.2}}


This specializes the elliptic-surface height and specialization results cited in the re-derivation. A complete re-derivation is given in Appendix~\ref{app:rederive:fam-src-D.1.2}.


\begin{theorem}[counterexample]\label{fam:src:D.1.3}\label{fam:st:D.1.3}
The proposed assertion that \emph{existence of a surjection onto the
specified \(E_t\) with elliptic kernel} is equivalent to \(6P(t)=O\)
is false. It remains true for the specified projection in \ref{fam:src:D.1.1}.
\end{theorem}
\Uses{the hypotheses, notation and definitions stated here; \cref{fam:src:D.1.1}; \cref{fam:src:C.2.1}}

\begin{proof}
Set \(t=32/5\). The exact fourth division polynomial is
\[
 \psi_4(P)=-\frac1{32}t^{10}(t-1)^4(5t-32),
\]
whereas \(\psi_2(P)=t^2(t-1)\ne0\). Thus \(P\) has exact order four.
On the rectangular branch write \(\tau=iT\), \(T>1\), with
\(j(iT)=1728\cdot32/5\). The chosen zero has real part \(1/2\), negative
imaginary part of absolute value less than \(T/2\), and order four;
therefore \(\mu=(2-\tau)/4\). Choose the concrete parameter
\[
 c=\left(\frac98-\frac{101}{4}\right)\tau-\frac94
       -\beta^{\rm part}(t).
\]
It is globally admissible. Indeed the positive-real Jacobi product at this
\((\mu,\tau)\) gives
\(\operatorname{Im}\beta^{\rm part}\geq-T-(3/\pi)\log3\).
Consequently
\(\operatorname{Im}c\leq-185T/8+(3/\pi)\log3<-22<-M\), by \ref{fam:src:C.2.1}.
This inequality proves admissibility without using the numerical estimate
\(c=-2.25-34.02109724218209472308\,i\).

The holomorphic linear map \(L:\mathbb C^2\to\mathbb C\),
\(L(\zeta)=3\zeta_1+4\zeta_2\), takes the four periods to
\[
 -101\tau,\qquad2\tau+2,\qquad3,\qquad4.
\]
These generate exactly \(\mathbb Z\tau+\mathbb Z\): the last two generate
\(1\), and the first two then generate \(\tau\), since
\(\gcd(101,2)=1\). Thus \(L\) descends to a surjective holomorphic group
homomorphism \(F_{32/5}\to E_{32/5}\) with connected one-dimensional
kernel. The kernel is compact as a closed subset of a compact torus,
hence an elliptic curve. Nevertheless \(6P=2P\ne O\).
The program \path{cd_torsion_independent.py} checks the displayed
division polynomials and four lattice images by exact arithmetic; its
command and retained output are in Appendix~\ref{app:I:programs}.
The admissibility inequality is proved analytically above.
\end{proof}

\begin{proposition}\label{fam:src:D.1.4}\label{fam:st:D.1.4}
For every smooth fibre and every admissible \(c\), its Neron--Severi lattice
is the following explicit subgroup of \(\mathbb Z^6\):
\begin{equation}\label{fam:eq:cd-NS}
 \left\{E:\ E_{12}-E_{13}\tau+(6E_{23}-E_{14})\mu+
 E_{24}\beta+E_{34}(6\mu^2-\tau\beta)=0\right\}.
\end{equation}
The complex orientation has intersection form
\(E^2=-2\operatorname{Pf}(E)\), where
\(\operatorname{Pf}(E)=E_{12}E_{34}-E_{13}E_{24}+E_{14}E_{23}\).
In particular the invariant class \(\eta=6\gamma\wedge\delta+u\wedge w\)
has \(\int_{F_t}\eta^2=-12\). Its Hermitian form is indefinite.
The fibres with lattice exactly \(\mathbb Z\eta\) have no compact curves
and algebraic dimension zero.
\end{proposition}
\Uses{the hypotheses, notation and definitions stated here; \cref{fam:src:C.3.1}; \cref{fam:src:D.1.6}}

\begin{proof}
On the universal cover the holomorphic two-form is the wedge of the two
rows of \(\Pi\). Its six coefficients in the integral coframe are
\(6\mu^2-\tau\beta,-\beta,6\mu,-\mu,\tau,1\).
An integral real two-form is of type \((1,1)\) exactly when its wedge
with this form vanishes: its \((0,2)\) part is then zero, and reality
also annihilates its \((2,0)\) part. Expanding the wedge gives
\eqref{fam:eq:cd-NS}. The exponential sequence identifies these integral
\((1,1)\) classes with Chern classes of holomorphic line bundles, proving
the lattice assertion.

The real period matrix in \ref{fam:src:C.3.1} has determinant \(TD<0\). Thus the
integral orientation \((\widehat\gamma,\widehat u,\widehat w,
\widehat\delta)\) is opposite to the complex orientation. Algebraically
\(E\wedge E=2\operatorname{Pf}(E)\gamma\wedge u\wedge w\wedge\delta\),
which proves the sign in the displayed intersection formula. In
particular the expression ``\(\eta^2=12\)'' in AC26, Lemma 9.4 is correct
as its displayed coefficient of \(\gamma\wedge u\wedge w\wedge\delta\),
and has the opposite sign as a complex-oriented numerical
self-intersection. We retain the coefficient identity and record the
orientation comparison explicitly.

For the precise Hermitian convention $H_E$ defined in \ref{fam:src:D.1.6}, direct
substitution of the four periods gives $H_\eta=-2G^{-1}$, where
\[
 G=\begin{pmatrix}2T&2U\\2U&\operatorname{Im}\beta/3\end{pmatrix},
 \qquad \det G=2TD/3<0.
\]
Thus $H_\eta$ has one positive and one negative eigenvalue, and every
nonzero multiple remains indefinite. The integration
current of an effective divisor, averaged under translations of the
compact torus, is a nonzero semipositive constant \((1,1)\)-form in its
Chern class. It cannot be an indefinite multiple of \(\eta\).
Thus no curve exists when \(\operatorname{NS}=\mathbb Z\eta\).
A nonconstant meromorphic function would have a nonempty polar divisor
(after subtracting a constant if necessary), so the algebraic dimension
is zero.
\end{proof}

\begin{proposition}\label{fam:src:D.1.5}\label{fam:st:D.1.5}
For a fixed \(c\), all but countably many smooth fibres have
\(\operatorname{NS}=\mathbb Z\eta\). Conversely, for every fixed smooth
\(t\), there are admissible parameters tending to the cusp of the
parameter disc for which \(F_t\) is an abelian surface, irrespective of
the order of \(P(t)\).
\end{proposition}
\Uses{the hypotheses, notation and definitions stated here; \cref{fam:eq:cd-NS}}

\begin{proof}
We verify the required independence of five holomorphic functions.
Suppose \(A+B\tau+C\mu+D\beta+E(6\mu^2-\tau\beta)=0\).
Apply \(g_0\), using \((\tau,\mu,\beta)\circ g_0=(\tau-1,\mu,\beta+1)\),
and subtract. The result is \(-B+D+E(\beta-\tau+1)=0\).
A second cusp difference gives \(2E=0\), hence \(E=0,D=B\).
Applying \(g_2\) to the remaining relation and multiplying the difference
by \(-\tau\) gives
\[
 6B\mu^2+C(\tau-1)\mu+B(\tau^2+3\tau+1)-C\tau=0.
\]
Its cusp difference gives \(C\mu+2B(\tau+1)-C=0\), whose next cusp
difference forces \(B=0\). Then \(C\mu=C\); the \(g_2\)-law forces
\(C=0\), and finally \(A=D=0\). Thus an integral class outside
\(\mathbb Z\eta\) has a nonidentically-zero holomorphic left side in
\eqref{fam:eq:cd-NS}; its zeros are discrete. Countability of \(\mathbb Z^6\)
proves the first assertion, with explicit verification of its reduction.

For the second assertion put
\(\beta_n=-3n^2\tau+12n\mu\),
\(c_n=\beta_n-\beta^{\rm part}(t)\). Its imaginary part tends to
minus infinity, so the already proved finite bound on \(M\) makes every
sufficiently large \(n\) admissible. The integral alternating matrix
with nonzero upper entries
\[
 E_{13}=-3n^2,\quad E_{14}=6n,\quad E_{23}=-n,\quad E_{24}=1
\]
satisfies \eqref{fam:eq:cd-NS} exactly at \(\beta_n\), and
\(\operatorname{Pf}(E)=-3n^2<0\). Its complex-oriented square is positive.
A nondegenerate Hermitian form in complex dimension two with positive
determinant is definite. Choosing its sign gives a positive integral
Riemann form, so \(F_t\) is projective. This supplies an explicit family
of projective fibres even at any chosen non-torsion value of the section.
\end{proof}

\begin{theorem}\label{fam:src:D.1.6}\label{fam:st:D.1.6}
Fix any admissible $c$. For each integral alternating matrix $E$, let
$Z_E(c)$ be the zero set on the marking cover of the explicit holomorphic
function in \eqref{fam:eq:cd-NS}, with $\beta=\beta^{\rm part}+c$ given by
\eqref{fam:eq:cd-beta}. Let $\pi$ be the cover map to $B^\circ$, and define
\[
 \mathcal E_c=\pi\!\left(\bigcup_{E\ne0,\ \operatorname{Pf}(E)=0}Z_E(c)\right),
 \qquad
 \mathcal A_c=\pi\!\left(\bigcup_{\operatorname{Pf}(E)<0}Z_E(c)\right).
\]
These give the exact loci of all elliptic quotients and all abelian
surfaces, respectively. The algebraic dimension of every smooth fibre is
\[
 a(F_t)=
 \begin{cases}
 2,&t\in\mathcal A_c,\\
 1,&t\in\mathcal E_c\setminus\mathcal A_c,\\
 0,&t\notin\mathcal E_c\cup\mathcal A_c.
 \end{cases}
\]
Together with \eqref{fam:eq:cd-NS}, these are an exact classification for every
$t$, by countable unions of explicitly given analytic zero sets; they do
not assert that the loci have a finite algebraic description. They include
all quotient maps, without requiring preservation of the first coordinate.
\end{theorem}
\Uses{the hypotheses, notation and definitions stated here; \cref{fam:eq:cd-NS}; \cref{fam:eq:cd-beta}; \cref{app:rederive:fam-src-D.1.6}}


This is the Appell--Humbert and Riemann criterion in the present lattice. A complete re-derivation is given in Appendix~\ref{app:rederive:fam-src-D.1.6}.


\begin{proposition}\label{fam:src:D.1.7}\label{fam:st:D.1.7}
There is an explicit nonempty locus of admissible members with ordinary
fibres of algebraic dimension one. On a sufficiently small punctured
cusp neighbourhood set
\[
 c(z)=-6\tau(z)-12\mu(z)-\beta^{\rm part}(z).
\]
Outside a countable set of $z$, the fibre of $X_{c(z)}$ over $\pi(z)$ has
Neron--Severi lattice $\mathbb Z\eta+\mathbb ZE$, where the nonzero upper
entries of $E$ are
$E_{13}=-6$, $E_{14}=-6$, $E_{23}=1$, $E_{24}=1$.
Its algebraic dimension is one. The section $P(\pi(z))$ can simultaneously
be required to be non-torsion.
\end{proposition}
\Uses{the hypotheses, notation and definitions stated here; \cref{fam:eq:cd-NS}; \cref{fam:src:D.1.6}; \cref{fam:src:C.2.1}; \cref{fam:src:D.1.2}}

\begin{proof}
At the specified parameter $\beta=-6\tau-12\mu$. Both $E$ and $\eta$
satisfy \eqref{fam:eq:cd-NS}, and exact multiplication gives
$\operatorname{Pf}(a\eta+bE)=6a^2$.
For any further integral class $F$, substitution in its period equation
gives a linear relation among $1,\tau,\mu,(\mu+\tau)^2$.
These four functions are linearly independent: applying $g_0$ to a
relation $A+B\tau+C\mu+D(\mu+\tau)^2=0$ and subtracting yields
$-B+D(1-2(\mu+\tau))=0$. A further cusp difference forces $D=0$,
and then $B=0$; nonconstancy of $\mu$ forces $A=C=0$.
Comparison of coefficients now gives
$F_{34}=F_{12}=0$, $F_{13}=-6F_{24}$, and
$6F_{23}-F_{14}=12F_{24}$, exactly
$F\in\mathbb Z\eta+\mathbb ZE$. Every other integral class therefore
has a discrete zero set along this period locus. Excluding their
countable union gives the asserted exact Neron--Severi lattice.
It contains the nonzero isotropic $E$ and no matrix of negative
Pfaffian; \ref{fam:src:D.1.6} gives algebraic dimension one.
Finally
$c=-5\tau-12\mu-(\beta^{\rm part}+\tau)$ has imaginary part tending
to minus infinity at the cusp. A sufficiently small cusp neighbourhood
is therefore admissible by \ref{fam:src:C.2.1}. The torsion locus of $P$ is countable
by \ref{fam:src:D.1.2}, so its exclusion still leaves such points in every open subset
of that neighbourhood. This supplies actual members and fibres, with
all their parameter functions specified.
\end{proof}

\subsection{Compact curves and surfaces}\label{fam:section:D.2}
\begin{proposition}\label{fam:src:D.2.1}\label{fam:st:D.2.1}
For a standard-twist member \(X_c\) with underlying manifold \(S^6\),
every compact irreducible complex curve lies in one fibre of \(f\), and
every compact irreducible complex surface is one of the reduced fibres.
\end{proposition}
\Uses{the hypotheses, notation and definitions stated here; the explicit constructions and calculations in the proof}

\begin{proof}
The given diffeomorphism with \(S^6\) implies \(H^2(X_c;\mathbb Z)=0\).
For a compact irreducible curve \(C\), let
\(\nu:\widetilde C\to C\hookrightarrow X_c\) be its normalization map.
The holomorphic map \(f\nu\) has
\(\deg (f\nu)^*\mathcal O_{\mathbb P^1}(1)=0\), since its Chern class is
pulled back from the zero class on \(X_c\). A nonconstant holomorphic map
of compact Riemann surfaces has positive degree: a regular value has a
positive number of preimages, all with positive local multiplicity.
Therefore \(f\nu\), and hence \(f|_C\), is constant.

An irreducible surface \(Z\) in the smooth threefold is an effective
Cartier divisor. If its image were all of \(B\), a general smooth fibre
would meet it in a nonempty effective divisor \(D\) whose line bundle
is \(\mathcal O_{X_c}(Z)|_{F_t}\). Its integral Chern class is zero.
Choose a flat K\"ahler form \(\omega\) on the torus \(F_t\). The equality
\(\int_{F_t}c_1(D)\wedge\omega=\int_D\omega\) has left side zero and right
side positive, a contradiction. Thus \(Z\) lies in a single fibre.
Each reduced fibre is irreducible of dimension two, so containment forces
equality. This uses the full map into \(X_c\), not a transfer of a
vanishing claim from its normalization.
\end{proof}

\begin{proposition}\label{fam:src:D.2.2}\label{fam:st:D.2.2}
There is no rational curve in \(X_c\setminus W_c\). There are elliptic
curves outside all three special fibres, in particular in each of the
two smooth fibres \(F_{t_\pm}\). All curves in \(W_c\) are precisely the
images of irreducible algebraic curves in \(\mathrm{dP}_6\) under its
finite normalization map. Distinct curves have the same image only for
the three paired boundary lines. All curves in a bielliptic fibre
\(S_j\) are images of irreducible curves in its finite product cover.
\end{proposition}
\Uses{the hypotheses, notation and definitions stated here; \cref{fam:src:D.2.1}; \cref{fam:src:D.1.1}; \cref{fam:src:C.4.1}}

\begin{proof}
By \ref{fam:src:D.2.1} every curve is vertical. A holomorphic map from
\(\mathbb P^1\) into a torus lifts to its vector-space universal cover,
since \(\mathbb P^1\) is simply connected, and its coordinate functions
are constant by compactness. Thus a smooth torus fibre contains no
rational curve. A map from \(\mathbb P^1\) into a bielliptic surface
lifts through its finite unramified torus cover by the same covering
criterion; the lifted map is constant. These arguments prove the first
assertion, including possibly singular rational images.
\ref{fam:src:D.1.1} gives an embedded connected elliptic kernel in each \(F_{t_\pm}\),
proving the second.

For the third assertion let \(\eta:\mathrm{dP}_6\to W_c\) be the finite
normalization morphism. Each component of the inverse image of a compact
curve is compact and one-dimensional and maps onto that curve, because
finite maps have zero-dimensional fibres. It is algebraic in the
projective surface \(\mathrm{dP}_6\). Conversely its image is a compact
analytic curve by properness. Away from the conductor, \(\eta\) is a
biholomorphism, so two distinct curves can have the same image only when
both lie in the conductor; its identifications are exactly the three
opposite pairs computed in \ref{fam:src:C.4.1}. This is also an explicit algebraic
parametrization: the curves of \(\mathrm{dP}_6\) are the three exceptional
lines and strict transforms of irreducible plane curves under the blow-up
of the three coordinate points. Their homogeneous equations and their
three multiplicities specify the strict transforms. The finite-product
cover statement for \(S_j\) follows from the identical proper-image and
finite-preimage argument applied to its unramified morphism
\(E_j\times F_j\to S_j\); curves in the product are algebraic.
\end{proof}


\paragraph{Cited input.} \label{fam:src:D.2.3}\label{fam:st:D.2.3}The source statement and its expanded derivation are retained in the proof compendium, Paper~I, label \nolinkurl{fam:src:D.2.3}.\par

\Uses{the hypotheses, notation and definitions stated here; \cref{fam:src:D.2.2}; \cref{fam:src:D.1.6}; \cref{fam:src:D.2.1}; \cref{app:rederive:fam-src-D.2.3}}


This is the effective line-bundle criterion for compact complex tori in the present lattice. A complete re-derivation is given in Appendix~\ref{app:rederive:fam-src-D.2.3}.




\subsection{The rank-three period quotient}\label{fam:section:C.3}
\begin{proposition}\label{fam:src:C.3.1}\label{fam:st:C.3.1}
At every point with \(D=0\), \(T>0\), the real-linear map
\(\Pi:\mathbb R^4\to\mathbb C^2\) has rank three. Write
\(\mu=a_0\tau+b_0\), \(a_0,b_0\in\mathbb R\), and put
\[
 k=(1,-6a_0,-6b_0,6a_0\operatorname{Re}\mu-\operatorname{Re}\beta).
\]
Its real kernel is \(\mathbb R k\). The subgroup \(\Pi\mathbb Z^4\) is
discrete exactly when \(k\in\mathbb Q^4\). In that case it is a
rank-three lattice in a real hyperplane and its quotient in \(\mathbb C^2\)
is a noncompact complex Lie group. In every other case the topological
quotient \(\mathbb C^2/\Pi\mathbb Z^4\) is non-Hausdorff.
\end{proposition}
\Uses{the hypotheses, notation and definitions stated here; the explicit constructions and calculations in the proof}

\begin{proof}
The real matrix, with output order
\((\operatorname{Re}\zeta_1,\operatorname{Im}\zeta_1,
\operatorname{Re}\zeta_2,\operatorname{Im}\zeta_2)\), is
\[
 \begin{pmatrix}6\operatorname{Re}\mu&\operatorname{Re}\tau&1&0\\
 6\operatorname{Im}\mu&T&0&0\\
 \operatorname{Re}\beta&\operatorname{Re}\mu&0&1\\
 \operatorname{Im}\beta&\operatorname{Im}\mu&0&0\end{pmatrix}.
\]
Its last three columns are real-linearly independent because \(T>0\).
Its determinant is \(TD\), so at zero determinant the rank is exactly
three. Substitution gives the displayed kernel, including its signs.
If the image subgroup is discrete, it is a lattice of rank at most three
in its real span. The homomorphism \(\mathbb Z^4\to\Pi\mathbb Z^4\)
therefore has a nonzero kernel, forcing the displayed kernel vector with
first coordinate one to be rational. Conversely a rational line has a
primitive integral generator extendable to a basis of \(\mathbb Z^4\).
The remaining three basis images are a real basis of the image hyperplane
and hence generate a discrete lattice. Its quotient leaves one unbounded
real direction. Finally a non-discrete finitely generated subgroup of
this real hyperplane cannot be closed: a closed subgroup is a product of
a real vector group and a discrete group, and a nonzero real vector group
is uncountable. A quotient topological group is Hausdorff exactly when its
identity subgroup is closed. This proves the last assertion.
\end{proof}

\begin{proposition}\label{fam:src:C.3.2}\label{fam:st:C.3.2}
The unique degenerate base point at the actual admissibility boundary is
\(t=1\). The boundary constants
\(c=x+i(3/\pi)\log6\) give
\[
 \ker_{\mathbb R}\Pi=\mathbb R(1,3,-3,-x-3/2).
\]
For rational $x+3/2=m/n$ in lowest terms, the quotient is a
finite holomorphic quotient of $E_i\times\mathbb C^*$, with a covering
of degree four. For irrational \(x+3/2\), the closure of the period subgroup is generated
by the real line \(\mathbb R(0,1)\) and the two periods
\((i,(1-i)/2)\), \((1,0)\). Its Hausdorff quotient is a real manifold
homeomorphic to \(\mathbb R\times(S^1)^2\).
\end{proposition}
\Uses{the hypotheses, notation and definitions stated here; \cref{fam:src:C.1.1}; \cref{fam:src:C.2.3}}

\begin{proof}
\ref{fam:src:C.1.1} gives \(\operatorname{Im}\beta=3/2\) at this value of \(c\).
The vector \(\widehat\gamma+3\widehat u-3\widehat w\) maps to
\((0,x+3/2)\), while \(\widehat\delta\) maps to \((0,1)\).
For $x+3/2=m/n$ in lowest terms, change coordinates by
$(\zeta_1,\zeta_2)\mapsto(\zeta_1,\zeta_2+\zeta_1/2)$.
The vectors $2\widehat u-\widehat\delta$ and
$2\widehat w-\widehat\delta$ then have periods $(2i,0)$ and $(2,0)$.
Bezout's identity applied to $\widehat\gamma+3\widehat u-3\widehat w$
and $\widehat\delta$ gives the period $(0,1/n)$.
These three periods generate a sublattice of index four in the rank-three
image: its two remaining generators are represented by
$(i,1/2)$ and $(1,1/2)$, each of order two modulo this sublattice.
The corresponding quotient is $E_i\times(\mathbb C/(n^{-1}\mathbb Z))$,
and $[z]\mapsto e^{2\pi inz}$ identifies its second factor with
$\mathbb C^*$. Lattice inclusion gives the claimed degree-four map.
For irrational $x+3/2$, its additive subgroup with $1$ is dense in the
indicated real line. The other two displayed periods are independent
modulo that line. This proves the closure and the quotient description.
By \ref{fam:src:C.2.3} this is the only degenerate torus cover at the admissibility
boundary; every other finite-base period matrix retains real rank four.
We have described the torus cover of the order-four filling, including
the actual quotient and its topology.
\end{proof}

\begin{remark}[Local boundary obstruction]\label{fam:src:C.3.3}\label{fam:st:C.3.3}
The local compactification obstruction at the actual boundary is localized at the order-four fibre over $t=1$. The proposed extension would require a compact complex-space replacement of the rank-three
quotient in \ref{fam:src:C.3.2}, over a neighbourhood of one boundary parameter,
with a proof that it extends the original cyclic filling and its overlap
with the ordinary torus family. Both rational and irrational values of
$x+3/2$ are included in the exact quotient description. The naive
quotients cannot be central fibres of a proper complex-space family:
for rational $x+3/2$ they are noncompact, and for irrational values
they are non-Hausdorff. \ref{fam:src:C.3.4} supplies a further obstruction to preserving
a holomorphic Weierstrass Albanese model in this parameter.
These prove failure of the stated specific replacements. The unique base-fibre location and exact local period obstruction are thereby determined. The general proper holomorphic extension is excluded by \cref{push:src:D.1.2,push:src:D.1.3}, which recovers the intrinsic divisor in a hypothetical extension.
\end{remark}
\Uses{the hypotheses, notation and definitions stated here; \cref{fam:src:C.3.2}; \cref{fam:src:C.3.4}; \cref{push:src:D.1.2,push:src:D.1.3}}


\begin{proposition}\label{fam:src:C.3.4}\label{fam:st:C.3.4}
On its half-plane of definition, the $j$-invariant of
$\operatorname{Alb}(S_2)$ has no meromorphic continuation across any
point of the line
$\operatorname{Im}c=(3/\pi)\log6$. In particular the rank-three quotient
at $t=1$ cannot be replaced by a Weierstrass family of elliptic Albanese
curves whose coefficients are holomorphic in the same parameter across
that line.
\end{proposition}
\Uses{the hypotheses, notation and definitions stated here; \cref{fam:section:A.4}; \cref{fam:src:C.1.1}}

\begin{proof}
The explicit Albanese lattice in \ref{fam:section:A.4} is
$\mathbb Z+(\kappa_2/2)\mathbb Z$, with
$\kappa_2=c+3/2-i(3/\pi)\log6$ by \ref{fam:src:C.1.1}. A positive-imaginary modulus is
$z=-\kappa_2/2$, so the invariant is $j(z)$. At a reduced rational
$r=a/b$ on the real boundary, choose an integral determinant-one modular
matrix with bottom row $(b,-a)$. Its value at $r+iy$ has imaginary part
$1/(b^2y)$. The convergent Fourier expansion $j=q^{-1}+744+O(q)$ therefore
gives growth of order $\exp(2\pi/(b^2y))$ as $y\downarrow0$.
A meromorphic function at $r$ is bounded by a fixed power of $1/y$,
so continuation there is impossible. A meromorphic continuation near an
irrational real point would also continue the same function near some
rational real point in that neighbourhood, which has just been excluded.
The affine change from $c$ to $z$ transfers this obstruction to the
stated line. A Weierstrass family with holomorphic coefficients has a
meromorphic $j$-invariant, expressed as the quotient of its cubic
invariants, so it cannot supply the indicated continuation. This
obstruction concerns the stated Albanese construction; it does not
silently assume that every possible compact threefold filling preserves
an Albanese model of its special surface.
\end{proof}

\subsection{Classification of the non-normal cusp surface}\label{fam:section:C.4}
\begin{theorem}\label{fam:src:C.4.1}\label{fam:st:C.4.1}
For the normalization of \ref{fam:src:C.1.1}, the three opposite-side identification
constants of \(W_c\), in the characters specified below, are
\[
 (g_1,g_2,g_3)=(-a^{-1},a^6,-a^{-7}e^{-2\pi ic}).
\]
Two such compact complex surfaces are biholomorphic exactly when
\(c'-c\in\mathbb Z\). The vertical \(\mathbb C^*\)-action inherited from
translation in the \(\widehat\delta\)-direction fixes exactly the double
curve represented by the pair of edges of direction \(e_2\).
\end{theorem}
\Uses{the hypotheses, notation and definitions stated here; \cref{fam:src:C.1.1}}

\begin{proof}
Keep the lattice of toric cocharacters \(\mathbb Z^2\), with
\(e_1=\widehat w\), \(e_2=\widehat\delta\), and cusp logarithm
\(s=(2\pi i)^{-1}\log t_c\). The cusp period matrix is
\([sB_0+C(t_c)\mid I]\), where
\[
 B_0=\begin{pmatrix}0&1\\-1&0\end{pmatrix},\qquad
 C(0)=\begin{pmatrix}6\mu_0&h_0\\b_0-h_0&\mu_0\end{pmatrix},
 \quad b_0=c.
\]
For \(v=e_1,e_2,e_1-e_2\), the gluing map from the side of direction
\(v\) to its opposite is the restriction of the toric action with lattice
label \(B_0v\). A character row \(r\) with \(rv=0\) pulls back by this
map to itself multiplied by
\(\exp(2\pi i\,rC(0)B_0v)\): the remaining power of \(t_c\) is
\(rB_0^2v=-rv=0\). Taking respectively
\(r=(0,1),(1,0),(1,1)\) gives exponents
\(-\mu_0,6\mu_0,-7\mu_0-c\). This proves the three constants, including
the cancellation of \(h_0\). These are morphisms of the whole projective
boundary lines, since multiplication by a nonzero constant extends over
both endpoints. The restrictions agree at both triple points, as the
ambient toric action supplies the same vertex identifications.

A biholomorphism of the two surfaces lifts uniquely to a biholomorphism
of their normalizations: integral closure of the reduced local rings is
functorial under an isomorphism. It preserves the inverse image of the
singular locus, the hexagon on \(\mathrm{dP}_6\). Its permutation of the
six components preserves their cyclic incidence, hence lies in the
order-twelve dihedral group. After composition with the corresponding
toric lattice automorphism, it fixes each boundary component. Each torus
character has its divisor supported on those components, so its pullback
divided by itself has divisor zero on the compact normalization and is
constant. Such an automorphism is therefore multiplication by an element
of \((\mathbb C^*)^2\). These torus multipliers cancel in the opposite-side
comparison because the same character is used on both sides.

The remaining six effective actions on the triple, with the central
lattice automorphism \(-I\) acting trivially, are described by the equations
\(g'_{k_j}{}^{\epsilon_j}=g_j\), where
\[
 ((k_1,\epsilon_1),(k_2,\epsilon_2),(k_3,\epsilon_3))
 \in\left\{
\begin{array}{lll}
((1,+),(2,+),(3,+)),&((1,-),(3,+),(2,+)),\\
((2,-),(3,-),(1,+)),&((2,+),(1,+),(3,-)),\\
((3,-),(2,-),(1,-)),&((3,+),(1,-),(2,-)).
\end{array}\right.
\]
This finite calculation is reproduced by \texttt{cd\_exact.py}, enumerating
all integral automorphisms of the six rays and transporting the characters.
Every nonidentity action requires one of
\(g_1^{-1}=g_1\), \(g_2^{-1}=g_2\), or equality of a positive real number
\(g_2^{\pm1}\) with a negative real number \(g_1^{\pm1}\).
All fail for \(a>1\). Therefore \(g'_3=g_3\), precisely
\(e^{-2\pi ic'}=e^{-2\pi ic}\). Conversely equality of the three constants
gives the identity gluing morphism on the normalization and hence a
biholomorphism of the quotients. This proves the unmarked assertion.

For the action of the subtorus with cocharacter \(e_2\), a one-dimensional
toric orbit on the normalization is pointwise fixed exactly when its
stabilizer line contains \(e_2\). Among the six rays this occurs exactly
for \(\pm e_2\). Their images form one double curve; all zero-dimensional
orbits map to its two endpoints. No dense two-dimensional orbit is fixed.
The quotient action has the same fixed set: a continuous orbit cannot
jump between the finitely many branches of a normalization fibre.
\end{proof}


\section{Real structures and their fixed three-spheres}\label{mapsI:real-forms}

We use the family marking
$\Lambda=\mathbb Z\langle\widehat\gamma,\widehat u,\widehat w,
\widehat\delta\rangle$ and its dual covector $\gamma$.
In this section $d=\widehat\delta$, $g=\widehat\gamma$,
$u=\widehat u$, and $w_{\Lambda}=\widehat w$; an unadorned $w$ in a
lattice expression denotes $w_{\Lambda}$, whereas the scalar parameter is
$w=\exp(-\pi i c)$. Put
\[
 e=\Pi_c d=(0,1)^t,\qquad E=d\otimes\gamma,
 \qquad \Pi_c=
 \begin{pmatrix}6\mu&\tau&1&0\\\beta^{\rm part}+c&\mu&0&1\end{pmatrix},
\]
\[
 D_c=\operatorname{Im}(\beta^{\rm part}+c)
       -6(\operatorname{Im}\mu)^2/\operatorname{Im}\tau<0.
\]
Thus $d$ is a lattice vector and $e$ is its image in the complex translation
space. The circle $\rho(s)$ is translation by $se$, $s\in\mathbb R/\mathbb Z$.
The clockwise root generators and the affine vectors are
\[
 A_1=\begin{pmatrix}1&0&0&0\\6&0&1&0\\-6&-1&-1&0\\-2&1&0&1\end{pmatrix},\quad
 A_2=\begin{pmatrix}1&0&0&0\\0&0&-1&0\\-6&1&0&0\\3&0&1&1\end{pmatrix},
\]
\[
 M_0=(A_1A_2)^{-1},\qquad
 v_1=(1,2,-4,0)^t,\qquad v_2=(-1,-3,3,0)^t.
\]
These are the same conventions as in the family construction. In cusp
coordinates $q=1/t$ and $s=(2\pi i)^{-1}\log q$, the bounded matrix $C_c(q)$
is defined by $Z_c=sB_0+C_c(q)$, with
$B_0=\left(\begin{smallmatrix}0&1\\-1&0\end{smallmatrix}\right)$.
In particular $\Pi_cg=(6\mu,-s+C_{c,21})^t$.
We write $D_1,D_2,D_{12}$ for the three cusp double curves belonging to
edges in directions $e_1,e_2,e_1-e_2$, respectively. Thus $D_2$ is the
fixed sphere of the circle action $\rho$, whereas $D_1$ is the double
curve along which the literal character has a limit. This distinction
is essential in the real calculation as well.

\begin{proposition}[Two conjugacy classes of real structures]\label{mapsI:src:A.3.1}
For a real member the antiholomorphic involutions have exactly two
conjugacy classes under holomorphic automorphisms. For any fixed
involutive splitting $\sigma$ they are represented by $\sigma$ and
$\rho(1/2)\sigma$. Every nonempty fixed set is a closed connected
orientable totally real three-manifold with the mod-two homology of
$S^3$. The circle $\rho$ preserves each such fixed set and acts
semifreely on it.
\end{proposition}
\Uses{the hypotheses, notation and definitions stated here; \cref{fam:src:B.4.2}; \cref{fam:src:B.4.1}; \cite[Theorem~3.1, p.~48]{Illman1972}; \cite[Theorem 4.1(1)]{Knudson2004}}


\begin{proof}
The real comparison of Proposition~\ref{fam:src:B.4.2} has derivative complex conjugation in the
second period coordinate. It therefore conjugates translation by
$b\in\C/\Z$ to translation by $\bar b$. In multiplicative notation its
action is $z\mapsto1/\bar z$. The square of $z\sigma$ is $z/\bar z$,
so it is an involution exactly for $z\in\R^*$. Conjugation by $y\in\C^*$
changes this coefficient to $|y|^2z$, giving exactly its two signs.
The automorphism computation used here is Proposition~\ref{fam:src:B.4.1}: the integral centralizer has linear part $aI$, the cusp and
finite norm equations force $a=1$ for a self-map, and the global
translation sheaf has $H^0=\C e$ with integral periods $\Z e$.
These are precisely the three steps needed to identify the full group
with $\C^*$, rather than only its identity component.

At a fixed point an antiholomorphic involution has a real three-dimensional
fixed tangent space, and its complexification is the full tangent space.
An invariant metric and its exponential map identify its fixed locus
locally with this real subspace. It is therefore a closed totally real
three-manifold. The smooth $\mathbb Z/2$ action admits an equivariant
triangulation by \cite[Theorem~3.1, p.~48]{Illman1972}; compactness and finite
dimension give the finite-dimensional $G$-CW hypothesis. Smith's theorem,
in the form of
\cite[Theorem 4.1(1)]{Knudson2004}, applies to the smooth finite action on
the sphere and gives the stated mod-two homology. A nonempty
three-dimensional fixed set is consequently connected, and its first
mod-two cohomology vanishes; its orientation class $w_1$ therefore
vanishes. On the unit circle $1/\bar z=z$, so $\sigma$ commutes with
$\rho$. Its stabilizers are inherited from the checked semifree action
on $X$ and hence are either trivial or the whole circle.
\end{proof}

We now fix a representative of each sign class, instead of leaving the
symbol $\sigma$ unspecified. Put
\[
 L_F=(2\pi i)^{-1}\log F,\qquad L_G=(2\pi i)^{-1}\log G,
 \qquad F^3=t,\quad G^4=t-1.
 \label{mapsI:real-eq-1}
\]
The roots are positive over $t>1$. We continue from that interval
through the upper half $t$-plane. Their logarithms are taken on the
orbifold uniformizing cover. Write $\iota$ for its reflection lift.
Every point of $-i\R_{>0}$ has a representative $c=1/2+iy$, and every
point of $i\R_{>0}$ has a representative $c=-1/2+iy$, with $y$ in the
admissible range. Set
\[
\begin{array}{c|c|c}
 &\alpha&S(t)\\\hline
 c=1/2+iy&1/3&2L_F-4L_G-1/6\\
 c=-1/2+iy&1/2&-2L_G+1/4
\end{array}
\qquad b_c(z)=\alpha\Pi_c(z)g+S(t(z))e.
\label{mapsI:real-eq-2}
\]

\begin{theorem}[Explicit real involutions and all fixed fibres]\label{mapsI:src:A.3.2}
On the ordinary period charts the formulas
\[
 \sigma_c[z,\zeta]=[\iota z,\bar\zeta+b_c(\iota z)],
 \qquad \sigma_c^- =\rho(1/2)\sigma_c
 \label{mapsI:real-eq-3}
\]
extend to antiholomorphic involutions of the compact threefold.
With these representatives their complete fixed-set data are as follows.
Here $2T^2$ means the disjoint union of two real two-tori; every displayed
circle in a double curve avoids its two triple-point endpoints.
\begin{center}\small
\begin{tabular}{cclccc}
\toprule
$w$&involution&image in $\RP^1$&on $S_1$&on $S_2$&on $W$\\
\midrule
$-ir$, $r>0$&$\sigma_c$&all $\RP^1$&$2T^2$&$2T^2$&$T^2\sqcup S^1_{D_2}\sqcup S^1_{D_{12}}$\\
$-ir$, $r>0$&$\sigma_c^-$&$[1,\infty]$&$\varnothing$&$T^2$&$S^1_{D_1}\sqcup S^1_{D_2}$\\
$ir$, $r>0$&$\sigma_c$&$\varnothing$&$\varnothing$&$\varnothing$&$\varnothing$\\
$ir$, $r>0$&$\sigma_c^-$&$[\infty,1]$ through $0$&$2T^2$&$T^2$&$S^1_{D_1}\sqcup S^1_{D_{12}}$\\
\bottomrule
\end{tabular}
\end{center}
The first interval is the closure of $t>1$; the second is the closure
of $t<1$. On every open real arc contained in the indicated image the
real regular fibre is $2T^2$. No triple point is real for these structures.
\end{theorem}
\Uses{the hypotheses, notation and definitions stated here; \cref{fam:src:B.4.2}; \cref{mapsI:real-eq-4}; \cref{mapsI:real-eq-2}; \cref{fam:src:A.3.1}; \cref{mapsI:real-eq-1}}


\begin{proof}
We first verify the global map. Use the lattice reflection from Proposition~\ref{fam:src:B.4.2},
\[
 H=\begin{pmatrix}1&0&0&0\\0&1&0&0\\-6&0&-1&0\\0&-1&0&-1\end{pmatrix},
 \quad E=d\otimes\gamma,
 \quad\Pi_{-\bar c}(\iota z)H=-\overline{\Pi_c(z)}.
\]
The initial map $\zeta\mapsto-\bar\zeta$ has target affine vectors
$-\varepsilon+d$ and $\varepsilon'-3d$. The normalization from this
target at parameter $c_0=-\bar c$ to the canonical datum at $c_0+1$
has linear part $-I$ and translation
\[
 b^{\rm norm}_{c_0}=\tfrac13\Pi_{c_0}g+(2L_F-4L_G)e.
 \label{mapsI:real-eq-4}
\]
Its period defects at the two finite generators are
$(0,-2,2,0)^t$ and $(0,0,2,0)^t$. To check them, the logarithmic increments
are $-2/3$ and $1$, and one multiplies the two displayed matrices $A_j$
by $g/3$. The finite difference vectors in the complex translation space are $2e$ and $-4e$, so subtracting
their logarithmic collars leaves holomorphic translations on both full
root discs. At the cusp
$2L_F-4L_G=s/3-(2\pi i)^{-1}\log(1-q)$ and
$\Pi g=(6\mu,-s+C_{21})^t$. The logarithmic terms cancel and give the
holomorphic toric unit
$(2\mu,C_{21}/3-(2\pi i)^{-1}\log(1-q))$.
Thus \eqref{mapsI:real-eq-4} extends over all three fillings, with its inverse.

For $c=1/2+iy$ compose this normalization with the initial reflection
and then with the constant translation $e/6$. This gives exactly the
first line of \eqref{mapsI:real-eq-2}. For $c=-1/2+iy$, additionally close the parameter
shift by the following explicit even-shift formula with $h=-2$,
\[
 b_{h,c_0}=-\frac h{12}\Pi_{c_0}g+
 \left(h(L_F-L_G)-\frac{h^2}{24}\right)e.
\]
A final constant translation $-e/4$ gives the second line of \eqref{mapsI:real-eq-2}.
The displayed translation has integral affine defects on the ordinary cover, holomorphic finite-chart corrections, and a nonvanishing toric-unit expression at the cusp; it represents the even shift of Proposition~\ref{fam:src:A.3.1}. Its cusp expression is checked immediately below. Indeed put $\ell=L_F-L_G=-s/12+\kappa_0(q)$ and
$C_{c,21}(q)=c+C_{0,21}(q)$, and set
$D_0(q)=C_{0,21}(q)-12\kappa_0(q)$.
The explicit even-shift translation becomes
\[
 b_{h,c}=
 \left(-h\mu/2,-h(c+D_0)/12-h^2/24\right)^t.
\]
The logarithmic terms have cancelled exactly. Exponentiating both
coordinates gives holomorphic nonzero torus multipliers on the whole cusp
disc; multiplication by these units extends over every toric stratum.
Compatibility with the period deck group is the integral identity
$\exp(2\pi i h\lambda_1)=1$ for $h\in2\mathbb Z$ and
$\lambda_1\in\mathbb Z$. The finite-chart defects are checked by the
same displayed $A_j$ and root logarithms used for the normalization.

Over $t>1$ we have
$\Re\mu=1/2$, $\Re\tau=\Re\beta^{\rm part}=0$ and the logarithms in
\eqref{mapsI:real-eq-1} are purely imaginary. The two squares are respectively the
integral period translations $2\widehat w$ and $3\widehat w$.
The square is a holomorphic automorphism and is consequently the
identity everywhere. This proves the involution claim.

We next give exact equations for all real fibres. Put $a=\Re c$ and
$R=-H+2aE$. In the chosen branch the reflection identifications on
$t>1$, $0<t<1$, $t<0$ are respectively $1,g_2,g_0^{-1}$.
They follow by reflecting the chosen upper-half-plane paths; equivalently
the root arguments are $(0,0)$, $(0,\pi/4)$ and $(\pi/3,\pi/4)$.
The three fibre involution matrices are
\[
 R_>=R,\qquad R_m=A_2^{-1}R,\qquad R_<=M_0R.
 \label{mapsI:real-eq-5}
\]
In marked real coordinates let $r_c(t)=\Pi_c(t)^{-1}b_c(t)$,
where $\Pi_c^{-1}$ is the real inverse. The fixed tori are exactly
\[
 \{x\in\R^4/\Lambda:(I-R_\bullet)x=
                       r_c(t)+\epsilon d/2\pmod\Lambda\},
 \qquad\epsilon=0,1.
 \label{mapsI:real-eq-6}
\]
This specifies the cosets, including their parameter dependence. Their
identity subgroups are $T_{\langle w,d\rangle}$ on both outer arcs and
$T_{\langle u+w,d\rangle}$ on the middle arc. When nonempty there are
two cosets. Relative to any one solution their offsets are
$(g+u)/2$ on the right arc and $g/2$ on the other two arcs.

For completeness, the finite arithmetic deciding existence can be
performed without approximating the periods. An affine involution
$x\mapsto R_\bullet x+\beta$ has a fixed point if and only if its
square translation $(I+R_\bullet)\beta$ belongs to
$(I+R_\bullet)\Lambda$. In \eqref{mapsI:real-eq-2} the square translation is
\[
 (I+R_\bullet)\alpha g+2\Re S\,d.
\]
For $c=1/2+iy$ the values $2\Re S$ on the three arcs are
$-1/3,-4/3,-2/3$; for $c=-1/2+iy$ they are $1/2,0,0$.
The resulting integral square vectors are respectively
$(2w,0,2w)$ and $(3w,d,3w)$. Adding $d/2$ to the involution adds $d$
to each square vector. Taking column Hermite normal forms of
$I+R_\bullet$ gives exactly the nonempty entries in the table.
All three $I-R_\bullet$ have nonzero Smith factors $1,2$.
This proves the coset counts and offsets; the complete exact computation
is in \texttt{code/maps\_audit/MAPS/code/real\_involutions.py}. In particular the lift on
the left arc is $g_0^{-1}$, not $g_0$; the latter would give a nonintegral
square vector and is the wrong continuation convention.

At $t=0$ the root degree is odd. Conjugation by the degree-three deck
group acts transitively on the three reflection lifts; the centralizer
of a reflection in that deck group is trivial.
Thus the quotient has exactly the same two real tori as the adjacent
nonempty real regular fibres. In the empty cases there are none. This
also follows by solving the limiting affine equations, and checks the
full multiplicity-three quotient rather than treating it as a reduced
unramified fibre.

At $t=1$ write the affine deck map as
$G(x)=A_2x+v_2/4$. The real root involution has matrix $R_>$ and
translation
\[
 b_-=(1/3,0,0,-1/6)^t,
 \qquad b_+=(1/2,0,0,1/4)^t
\]
for the two parameter signs. The real locus of the quotient is the
quotient of the fixed sets of $G^kh$, $0\leq k<4$. Its even and odd
lift parities are identified internally by $G$, while $G^2$ preserves
each parity. Direct integer congruences show that the nonempty lift
parities are: both for negative $w$ and $\epsilon=0$; only even for
negative $w$ and $\epsilon=1$; neither for positive $w$ and $\epsilon=0$;
only odd for positive $w$ and $\epsilon=1$. Every nonempty lift has
two tori, exchanged by $G^2$. Each parity therefore contributes one
torus downstairs, not a Klein bottle. The exhaustive congruence check
\texttt{code/maps\_audit/MAPS/code/real\_finite\_four.py} uses denominator $24$: translation
denominators divide $12$, and the nonzero Smith factors are $1,2$, so
every component has such a representative. It checks the component
permutation, not just the dimension of the fixed eigenspace. The normal
real line has fourth power of sign positive for the even parity and
negative for the odd parity. This establishes the two one-sided
endpoints at $t=1$ recorded in the table.

It remains to compute the cusp, which also excludes hidden real
components confined to a singular fibre. Put
$A=C_{12}(0)$; the cusp normalization gives
\[
 \Re\mu(0)=1/2,\quad A\in i\R,\quad C_{21}(0)=c-A,
 \quad C_{11}(0)=6\mu(0),\quad C_{22}(0)=\mu(0).
\]
In a toric lift the involution is
$x\mapsto U/\bar x$, with
\[
\begin{array}{c|cc}
 &b_1(0)&b_2(0)\\\hline
 w=-ir&2\mu(0)&C_{21}(0)/3-1/6\\
 w=ir&3\mu(0)&C_{21}(0)/2+1/4
\end{array}
\qquad U_j=e^{2\pi i b_j(0)}.
\label{mapsI:real-eq-7}
\]
Thus $U_2>0$ in both rows, whereas $U_1>0$ in the first and $U_1<0$
in the second. Replacing $\sigma_c$ by $\sigma_c^-$ negates $U_2$.
The smooth central orbit has a real torus precisely when both units
are positive, with equations $|x_j|^2=U_j$.

For an edge $v$, set $n=(-v_2,v_1)^t$. The period translate bringing
the inverted edge back is $B_0n=v$. Its curve coordinate $z=x^n$ has
involution
\[
 z\longmapsto K_v/\bar z,
 \qquad K_v=U^n e^{2\pi i n^tC(0)n}.
 \label{mapsI:real-eq-8}
\]
The signs for $(D_1,D_2,D_{12})$ are $(-,+,+)$ for negative $w$
and $(-,-,-)$ for positive $w$. Negating $U_2$ reverses the first and
third signs, but not the second. A positive sign gives the single
circle $|z|^2=K_v$; a negative sign gives no fixed point. Central
inversion exchanges the two triangle classes of the periodic fan, so
it exchanges the two triple points. These calculations prove every
cusp entry. Finally a nonempty fixed three-manifold cannot be confined
to the special fibres: a holomorphic function constant on an open
piece of a totally real submanifold is locally constant. This also
proves that the positive-$w$ representative with no regular real
fibres is entirely fixed-point free.
\end{proof}

\begin{theorem}[The real three-manifolds]\label{mapsI:src:A.3.3}
Every nonempty fixed set in Theorem~\ref{mapsI:src:A.3.2} is diffeomorphic to $S^3$.
Its fundamental group is trivial and its integral homology is
$(\Z,0,0,\Z)$. For negative imaginary $w$, the restricted circle action
is the standard semifree action on $S^3$ with fixed circle
$D_2^{\sigma_c}$. For positive imaginary $w$, the nonempty real form
has a free circle action isomorphic, up to circle and orientation
conventions, to the Hopf action.
\end{theorem}
\Uses{the hypotheses, notation and definitions stated here; \cref{mapsI:src:A.3.2}}


\begin{proof}
For negative imaginary $w$ both fixed manifolds meet $D_2$ in one circle.
Their circle action is semifree, so the orbit space is a compact
orientable surface with boundary. The local quotient at a fixed circle
is the quotient of $S^1\times\C$ by scalar rotation on $\C$, hence is
a boundary collar. Over the surface interior the principal circle bundle
is trivial because the surface has boundary and $H^2=0$. Changes of
boundary trivialization extend over the attached normal discs by
rotation. Thus the three-manifold is
\[
 (B\times S^1)\ \cup_{\partial B\times S^1}\ (\partial B\times D^2).
\]
If $B$ has genus $g_B$ and $b_B$ boundary components, van Kampen or its
surface handle decomposition gives first homology
$\Z^{2g_B+b_B-1}$. Smith's theorem makes this rank zero. Since $b_B>0$,
we have $g_B=0,b_B=1$. The displayed manifold is then
$\partial(D^2\times D^2)=S^3$, with the indicated circle action. This
derives the smooth recognition from an explicit surface decomposition.

For positive imaginary $w$ the nonempty real manifold is free under
$\rho$. The two real tori on $t<1$ join at the single real torus over
$1$. At $0$ the odd root comparison is a torus product. At infinity
they are capped by the real neighbourhoods of $D_1$ and $D_{12}$.
These are solid tori: the primitive circle in the edge direction
collapses and the remaining real circle is their core. In the outer
real-fibre lattice $\Z\langle w,d\rangle$ the two cap meridians are
$w$ and $w-d$.

We check the intervening torus identification. Put
$B_1=A_2^{-1}A_1^{-1}A_2$ and $C=B_1^2$. The two real-arc involutions
satisfy $R_<R_m=B_1$, $R_mB_1R_m=B_1^{-1}$ and $B_1^3=I$.
Moving between the real root lifts requires
$B_1^kR_mB_1^{-k}=B_1R_m$, so $2k\equiv1\pmod3$ and $k=2$.
Thus the actual odd-root transport is $C$, not merely an arbitrary
matrix conjugating the two reflections. It satisfies $R_<=CR_mC^{-1}$, and
$C(u+w)=-w+d$. The exchange across the multiplicity-four real torus
is $A_2^2$; after this odd-root transport it is
$CA_2^2C^{-1}$, whose restriction to $\langle w,d\rangle$ is
$w\mapsto-w$, $d\mapsto d$. Consequently the meridians, compared on
one boundary torus, have determinant of absolute value one. The
resulting genus-one Heegaard splitting is $S^3$: primitive meridians
intersecting once can be taken to the two coordinate circles by an
integral torus automorphism, giving the usual two solid tori of $S^3$.
The free circle $d$ then has Euler number of absolute value one, so
the displayed bundle over $S^2$ is the Hopf bundle. This proves the
integral assertions, beyond the mod-two information.
\end{proof}

\begin{proposition}[Restrictions of the repaired circle quotient]\label{mapsI:src:A.3.4}
The repaired map $e_{\rm q}$ of Theorem~\ref{mapsI:src:A.2.1} may be chosen $\rho$-invariant. Its restriction to
every nonempty real form is null-homotopic. On an ordinary real fibre
its image consists of the two quotient-circle values
\[
 \theta=\tfrac12\gamma(r_c(t)),\qquad
 \theta=\tfrac12\gamma(r_c(t))+\tfrac12\pmod\Z,
 \quad
 \gamma(r_c(t))=\alpha+\frac{\Im S(t)}{D_c(t)},
 \label{mapsI:real-eq-9}
\]
with $\alpha,S$ from \eqref{mapsI:real-eq-2}. The circle action is along each real
two-torus and is annihilated by $e_{\rm q}$.
\end{proposition}
\Uses{the hypotheses, notation and definitions stated here; \cref{mapsI:src:A.2.1}; \cref{mapsI:real-eq-2}}


\begin{proof}
The literal quotient is circle invariant. In a cusp extension
$\theta':N_0\to S^1$, each orbit circle has degree zero: its class is
$d$, which vanishes in $\pi_1(N_0)$. The quotients
$\theta'(\rho(s)x)/\theta'(x)$ therefore have a continuous real lift
vanishing at $s=0$ and periodic in $s$. Averaging this real lift over
$s\in\R/\Z$ changes $\theta'$ to a circle-invariant map. The cocycle
identity proves invariance, and on an invariant outer collar the lift
is zero, so the replacement agrees exactly with the prescribed map.
The same construction is smooth. Thus $e_{\rm q}|_{\operatorname{Fix}\sigma}$
factors through its circle orbit space, which is $D^2$ or $S^2$ by the
preceding proof. Both have only two-dimensional homotopy type and
$S^3$ is two-connected, so this restriction is null-homotopic.
The first row of $I-R_\bullet$ is $(2,0,0,0)$, giving the two values
in \eqref{mapsI:real-eq-9}. The real inverse of the period matrix sends $S e$ to
first coordinate $\Im S/D_c$; this follows by subtracting
$(\Im\mu/\Im\tau)$ times its first imaginary-coordinate equation
from the second. Both real-fibre identity subgroups have zero
$\gamma$-coordinate, proving the remaining assertions.
\end{proof}

The formulas specify the embedded fixed sets by their torus cosets and
their actual finite and toric attachments. No ambient-isotopy
classification or value of the Haefliger knot invariant of these
embedded three-spheres is asserted here.



\paragraph{Cited input.} \label{push:src:C.3.5}\label{mapsI:closing:eq:C-explicit-closing}\label{mapsI:closing:eq:C-twelve-scalar}The source statement and its expanded derivation are retained in the proof compendium, Paper~I, label \nolinkurl{push:src:C.3.5}.\par

\Uses{the hypotheses, notation and definitions stated here; \cref{fam:src:A.3.1}; \citet[Proposition~7.8 and Appendix~C, Proposition~C.2]{DuEtAl2026v1}; \cref{app:rederive:push-src-C.3.5}}


The raw even-shift formula and its coherent correction are due to \citet[Proposition~7.8 and Appendix~C, Proposition~C.2]{DuEtAl2026v1}. A complete re-derivation is given in Appendix~\ref{app:rederive:push-src-C.3.5}.



